\documentclass[UTF8,12pt,a4paper,fontset=fandol]{ctexart}

% 支持从项目根目录或当前笔记目录编译。
\makeatletter
\def\input@path{{../}{../../}}
\makeatother
\usepackage{researchnotes}

\title{拉普拉斯定理笔记}
\author{}
\date{}

\begin{document}
\maketitle

\section{从按一行展开到按多行展开}

行列式按一行展开，将一个 $n$ 阶行列式写成若干个元素与其代数余子式的乘积之和。
如果矩阵的零元素集中在若干行和若干列组成的区域内，逐行展开可能使这种整体结构不易看清。
\keyterm{拉普拉斯展开定理}（Laplace expansion theorem）允许固定若干行一起展开，
把一个高阶行列式写成若干个低阶行列式乘积之和。按一行或一列展开是它的特殊情形。

这一方法并非一般数值行列式的首选算法：展开项较多时，计算量可能很大。
它的主要用途是利用成组的零元素简化计算，以及证明行列式的结构性公式。
下文先说明子式与符号规则，再给出定理和证明，最后通过例题说明其作用及适用范围。
全文以实数域或复数域上的方阵为讨论对象。

\section{子式、余子式与符号规则}

设 $A=(a_{ij})$ 是 $n$ 阶方阵，$n\geq2$，并记 $[n]=\{1,\ldots,n\}$。
取整数 $1\leq k<n$，以及两个各含 $k$ 个指标的集合
\[
  I=\{i_1<\cdots<i_k\},\qquad J=\{j_1<\cdots<j_k\}.
\]
记 $I^c=[n]\setminus I$、$J^c=[n]\setminus J$，并记
\[
  s(I)=i_1+\cdots+i_k,\qquad s(J)=j_1+\cdots+j_k.
\]
用 $A[I,J]$ 表示保留 $I$ 中各行与 $J$ 中各列得到的子矩阵；
所有保留的行、列均按原来的先后顺序排列。

\begin{definition}[子式与余子式]
\label{def:laplace-minors}
行列式
\[
  M_{I,J}=\det A[I,J]
\]
称为 $\det A$ 的一个 $k$ 阶\keyterm{子式}（minor）。
删去 $I$ 中各行和 $J$ 中各列，剩下的行列式
\[
  M^{\mathrm{c}}_{I,J}=\det A[I^c,J^c]
\]
称为该子式的\keyterm{余子式}（complementary minor）。
相应的\keyterm{代数余子式}（signed complementary minor）定义为
\begin{equation}
  C_{I,J}=(-1)^{s(I)+s(J)}M^{\mathrm{c}}_{I,J}.
  \label{eq:laplace-cofactor}
\end{equation}
\end{definition}

这里的余子式为 $n-k$ 阶，而非固定为 $n-1$ 阶。
当 $k=1$ 时，这一定义就是元素 $a_{ij}$ 的通常代数余子式（cofactor）。
计算时需要区分三部分：展开符号 $(-1)^{s(I)+s(J)}$、子式自身的值、余子式自身的值。
展开符号为正，并不意味着这一项的数值为正。

\section{拉普拉斯展开定理及其证明}

\begin{theorem}[按固定的多行展开]
\label{thm:laplace-row}
设 $A$ 是 $n$ 阶方阵，$n\geq2$，$1\leq k<n$。
任意固定一个含 $k$ 个指标的行集合 $I\subseteq[n]$，则
\begin{equation}
  \det A
  =\sum_{\substack{J\subseteq[n]\\|J|=k}}
  (-1)^{s(I)+s(J)}\det A[I,J]\det A[I^c,J^c]
  =\sum_{\substack{J\subseteq[n]\\|J|=k}}M_{I,J}C_{I,J}.
  \label{eq:laplace-row}
\end{equation}
即固定的 $k$ 行所构成的所有 $k$ 阶子式，各自乘以其代数余子式，再相加，得到原行列式。
\end{theorem}

定理中的关键是\keyterm{固定行集合，遍历列集合}。
未经零项筛选时共有 $\binom nk$ 个展开项；如果同时遍历所有行集合与列集合，
就不再是式~\eqref{eq:laplace-row}，而会重复计数。

\begin{proof}
证明的目标是把行列式定义中的排列项分组，并说明每一组为何形成两个行列式的乘积。
记 $S_n$ 为 $[n]$ 上所有排列的集合，$\operatorname{sgn}(\sigma)$ 为排列 $\sigma$ 的符号。
由行列式的排列定义，
\[
  \det A=\sum_{\sigma\in S_n}\operatorname{sgn}(\sigma)
  \prod_{i=1}^n a_{i,\sigma(i)}.
\]
每个排列 $\sigma$ 把固定的行集合 $I$ 对应到唯一的列集合 $J=\sigma(I)$。
因此，可按 $J$ 将全部排列项分成互不相交的组。

固定其中一组 $J$。把 $I$ 中的行按原顺序移至最前面，其余行保持相对顺序；
把 $J$ 中的列也作相同处理。移动行所需的相邻交换次数为
\[
  \sum_{p=1}^k(i_p-p)=s(I)-\frac{k(k+1)}2,
\]
因为第 $i_p$ 行之前恰有 $i_p-p$ 个未被选中的行需要越过。
同理，列交换次数为 $s(J)-k(k+1)/2$。
故行、列重排合计产生的符号为
\[
  (-1)^{s(I)+s(J)-k(k+1)}=(-1)^{s(I)+s(J)},
\]
其中用到了 $k(k+1)$ 为偶数。

重排后的矩阵具有形式
\[
  \begin{pmatrix}
    A[I,J]&A[I,J^c]\\
    A[I^c,J]&A[I^c,J^c]
  \end{pmatrix}.
\]
本组中的排列满足 $\sigma(I)=J$，所以在重排后，它们将前 $k$ 行对应到前 $k$ 列，
并将其余行对应到其余列。因此，本组的每个乘积项只使用两个对角块中的元素。
这类排列由前 $k$ 个指标上的一个排列和后 $n-k$ 个指标上的一个排列唯一确定；
两部分之间没有逆序，故符号是两个排列符号的乘积。
分别对两部分排列求和，按行列式定义恰好得到
\[
  \det A[I,J]\det A[I^c,J^c].
\]
还原行、列顺序时再乘以 $(-1)^{s(I)+s(J)}$，便得到本组在原行列式中的贡献。
对所有含 $k$ 个指标的 $J$ 求和，即得式~\eqref{eq:laplace-row}。
\end{proof}

\begin{corollary}[按固定的多列展开]
在定理~\ref{thm:laplace-row} 的矩阵与阶数条件下，任意固定含 $k$ 个指标的列集合 $J$，则
\begin{equation}
  \det A=\sum_{\substack{I\subseteq[n]\\|I|=k}}
  (-1)^{s(I)+s(J)}\det A[I,J]\det A[I^c,J^c].
  \label{eq:laplace-column}
\end{equation}
\end{corollary}

\begin{proof}
将定理~\ref{thm:laplace-row} 应用于转置矩阵 $A^{\mathsf T}$，并利用行列式在转置下不变，
即得到式~\eqref{eq:laplace-column}。
\end{proof}

当 $I=\{i\}$ 时，式~\eqref{eq:laplace-row} 化为熟悉的按第 $i$ 行展开公式
\[
  \det A=\sum_{j=1}^n a_{ij}(-1)^{i+j}
  \det A[\{i\}^c,\{j\}^c].
\]
因此，多行展开与一行展开遵循同一个原理，只是分组方式不同。

\section{计算例题：只保留两个展开项}

\begin{example}
\label{ex:laplace-two-terms}
计算四阶行列式
\[
  D=\begin{vmatrix}
    1&0&2&0\\
    0&1&0&0\\
    3&4&5&6\\
    7&8&9&10
  \end{vmatrix}.
\]
\end{example}

\begin{solve}
固定前两行，即 $I=\{1,2\}$，因此 $s(I)=3$。
从四列中选两列原本有六种组合，但不选第 $2$ 列时，子式的第二行全为零；
选第 $4$ 列时，子式又含有零列。因此，只有 $J=\{1,2\}$ 和 $J=\{2,3\}$ 可能贡献非零项。

当 $J=\{1,2\}$ 时，展开符号为 $(-1)^{1+2+1+2}=1$，对应项为
\[
  \begin{vmatrix}1&0\\0&1\end{vmatrix}
  \begin{vmatrix}5&6\\9&10\end{vmatrix}
  =1\cdot(50-54)=-4.
\]
当 $J=\{2,3\}$ 时，展开符号为 $(-1)^{1+2+2+3}=1$，对应项为
\[
  \begin{vmatrix}0&2\\1&0\end{vmatrix}
  \begin{vmatrix}3&6\\7&10\end{vmatrix}
  =(-2)\cdot(30-42)=24.
\]
根据式~\eqref{eq:laplace-row}，两项相加可得
\[
  D=-4+24=20.
\]
第二项的展开符号为正，但子式与余子式的值均为负，计算时不能把这些符号混为一谈。
\end{solve}

本例适合熟悉多行展开的步骤，却不能说明它总比一行展开方便。
原行列式第二行只有一个非零元素，直接按第二行展开同样简洁：
\[
  D=\begin{vmatrix}1&2&0\\3&5&6\\7&9&10\end{vmatrix}
  =(50-54)-2(30-42)=20.
\]
选用哪种展开方式，应根据零元素的分布决定，而不必刻意使用多行展开。

\section{结构性应用：分块三角行列式}

\begin{proposition}[分块三角矩阵的行列式]
\label{prop:laplace-block-triangular}
设 $B$ 是 $r$ 阶方阵，$C$ 是 $s$ 阶方阵，$r,s\geq1$，
$E$ 是 $s\times r$ 矩阵，$F$ 是 $r\times s$ 矩阵。则
\begin{equation}
  \det\begin{pmatrix}B&0\\E&C\end{pmatrix}
  =\det B\det C,
  \qquad
  \det\begin{pmatrix}B&F\\0&C\end{pmatrix}
  =\det B\det C.
  \label{eq:laplace-block-triangular}
\end{equation}
这里的零矩阵分别为 $r\times s$ 与 $s\times r$；不要求 $B$ 或 $C$ 可逆。
\end{proposition}

\begin{proof}
对左侧的分块下三角矩阵，固定前 $r$ 行进行拉普拉斯展开。
只要所选的 $r$ 列包含后 $s$ 列中的任何一列，对应的 $r$ 阶子式就含零列，因而为零。
所以，只有选取前 $r$ 列这一项可能非零。
此时子式为 $\det B$，余子式为 $\det C$，展开符号为
\[
  (-1)^{2(1+\cdots+r)}=1.
\]
这就证明了第一个等式。对分块上三角矩阵转置，再应用第一个等式及转置不改变行列式的性质，
即可得到第二个等式。
\end{proof}

\begin{example}[含任意参数的行列式]
设 $a,b,c,d$ 为任意实数，计算
\[
  D=\begin{vmatrix}
    1&2&0&0\\
    3&4&0&0\\
    a&b&5&6\\
    c&d&7&8
  \end{vmatrix}.
\]
\end{example}

\begin{solve}
右上角的 $2\times2$ 子矩阵全为零，故由命题~\ref{prop:laplace-block-triangular}，
\[
  D=\begin{vmatrix}1&2\\3&4\end{vmatrix}
    \begin{vmatrix}5&6\\7&8\end{vmatrix}
   =(-2)(-2)=4.
\]
这不仅简化了运算，还说明左下角的参数 $a,b,c,d$ 完全不影响该行列式的值。
对于更高阶的同类矩阵，也可以直接将问题拆成两个较低阶行列式的计算。
\end{solve}

\begin{remark}[零块条件不能随意去掉]
式~\eqref{eq:laplace-block-triangular} 依赖分块三角结构，并不是任意分块矩阵都满足
\[
  \det\begin{pmatrix}B&F\\E&C\end{pmatrix}=\det B\det C.
\]
例如，将各块取为一阶矩阵，便有
\[
  \begin{vmatrix}1&1\\1&1\end{vmatrix}=0\neq1\cdot1.
\]
当两个非对角块均可能非零时，不能仅计算两个对角块的行列式之积。
\end{remark}

\section{如何判断是否适合使用}

拉普拉斯展开提供的是一个恒等式，恒等式本身并不保证直接计算时高效。
固定 $k$ 行后共有 $\binom nk$ 个候选项，每一项又包含两个行列式。
如果大部分子式都非零，直接展开常会增加工作量。

实际计算时，可以先检查矩阵的结构，再选择方法：
\begin{itemize}
  \item 某一行或某一列只有少数非零元素时，优先考虑按该行或该列展开。
  \item 若干行中的非零元素集中在少数列内时，考虑按这些行一起展开，并先排除零子式。
  \item 出现分块三角结构时，直接使用式~\eqref{eq:laplace-block-triangular}。
  \item 缺少明显稀疏结构的一般数值矩阵，通常采用初等行变换化为三角矩阵，
  并记录交换行和行倍乘对行列式的影响。
\end{itemize}

在每次标量算术运算计为常数代价的模型下，用通常的高斯消元计算一般 $n$ 阶行列式需要
$O(n^3)$ 次算术运算；若遇到零主元，需要交换行，若剩余主元列无法找到非零元素，则可判定行列式为零。
这一运算次数的比较不涉及浮点舍入误差，也不计精确运算中数值位数增长的代价。
它说明拉普拉斯展开的价值主要在于利用结构和推导公式，而不是作为所有行列式的通用快速算法。

\end{document}
